\scanpage{104}
\margin{\unsure{!}}

\diagram{$\begin{array}{c@{\quad}ccc}
 &\multicolumn{2}{c}{\overbrace{\hphantom{\text{poids 0}\qquad\text{poids 1}}}^{\text{\deleted{\unsure{pour}} \quad cas propre, ou du }H^1_!}}&\\
 H^1&\text{poids }0&\text{poids }1&\text{poids }2\\
 &\text{niveau }0&\text{niveau }1&\text{niveau }0\\
 &&\multicolumn{2}{c}{\underbrace{\hphantom{\text{niveau 1}\qquad\text{niveau 0}}}_{\text{cas lisse}}}
\end{array}$}
\margin{morceaux purs sont semi-simples}

\diagram{$\begin{array}{c@{\quad}c|ccccc}
 \multicolumn{2}{c}{}&\multicolumn{3}{c}{\overbrace{\hphantom{0\qquad1\qquad2}}^{\text{cas propre}}}&&\\
 H^2_{\deleted{\unsure{!}}}&\uline{\text{poids}}&0&1&2&3&4\\
 &\text{niveau}&0&1&2&1&0\\
 \multicolumn{2}{c}{}&&&\multicolumn{3}{c}{\underbrace{\hphantom{2\qquad1\qquad0}}_{\text{cas lisse}}}
\end{array}$}

morceaux purs de poids $\ne 2$ sont semi-simples.

$\Vert$ Le morceau pur de poids $2$ est une extension par un motif de Tate $Q(-1)^{r}$.

$H^2_!$ \quad il semble qu'il puisse y avoir mélange de poids dans le mauvais \unsure{rdre|ordre|sens}, à cause d'une extension d'un motif \unsure{effectif|affectif} de poids $0$ par un autre de poids $1$ ; en plus de cela, il apparaît une extension, peut-être non triviale, d'un motif effectif de poids $0$ \inserted{resp. $1$} par un autre.
% archive sheet number: 103
